\chapter{Discussion and Conclusion}

\vspace*{0.3cm}

The results show a clear tradeoff between model quality and computational
efficiency. Adam is the best optimizer for reducing training and validation loss,
which makes it the strongest quality baseline. However, it is slower than
SGD-based methods and has higher inference latency. For Edge AI, the best
optimizer is not necessarily the one with the lowest loss, because deployment
also depends on latency, throughput, memory usage, and model size.

SGD gives the best efficiency among the trained optimizers, while SGD with
momentum improves the validation loss with only a moderate increase in cost.
AdamW with clipping and scheduling is also a useful compromise because it gives
better quality than SGD while remaining more efficient than Adam. L-BFGS and
Newton-CG are important from an optimization theory perspective, but their high
memory and computational costs make them less practical for transformer
fine-tuning in an Edge AI context.

The project has several limitations. The dataset is intentionally small, with
2000 training examples and 500 validation/test examples. The benchmark is also
short, using only 10 training steps and a batch size of 4, so it compares early
optimizer behavior rather than complete convergence. The results also depend on
the GPU used during execution, and energy consumption could not be fully analyzed
without reliable power telemetry.

Overall, the project demonstrates that optimizer selection must be guided by the
target objective. Adam should be selected when fine-tuning quality is the main
priority. SGD with momentum or AdamW are better candidates when the objective is
to balance quality and efficiency for Edge AI deployment. The results provide a
baseline for future compression techniques such as quantization, pruning, and
knowledge distillation, which can further reduce memory usage and latency while
preserving acceptable model quality.
